\documentclass[11pt]{article}
\usepackage{amsmath,amssymb,amsthm}
\usepackage[margin=1in]{geometry}
\newtheorem{theorem}{Theorem}
\newcommand{\Tr}{\operatorname{Tr}}
\begin{document}
\title{An even-order pure-state equality obstruction\\and the exclusion of $N=16$, $t=6$}
\author{Catalogue research note}
\date{31 October 2026}
\maketitle

Write $V_j$ for the spin-$j$ representation, of dimension $2j+1$.
A pure $t$-anticoherent (AC) state in $V_j$, $j=N/2$, has vanishing
expectations of all irreducible spin tensors of ranks $1,\ldots,t$.
Consequently, if an operator $A$ has only tensor ranks at most $t$, then
\[
 \langle\psi|A|\psi\rangle=\frac{\Tr A}{N+1}.
\]
We use Hilbert--Schmidt orthonormal tensors throughout.

\begin{theorem}
There is no pure sixth-order AC state on sixteen symmetric qubits.
In particular, the maximal sixth-order AC-subspace dimension is
$K(16,6)=0$.
\end{theorem}

\begin{proof}
Suppose $|\psi\rangle\in V_8$ were $6$-AC.  Consider the rectangular
tensor operators $E_{\ell q}:V_8\to V_7$, for $\ell=1,2,3$ and
$q=-\ell,\ldots,\ell$.  There are $3+5+7=15$ such operators, and
$\dim V_7=15$.  Their explicit matrix elements are
\[
 (E_{\ell q})_{a,m}=(-1)^{7-a}\sqrt{2\ell+1}
 \begin{pmatrix}7&\ell&8\\-a&q&m\end{pmatrix},
 \qquad -7\le a\le7,\quad -8\le m\le8.
\]
Orthogonality of the Wigner $3j$ symbols gives
$\Tr(E_{\ell q}^\dagger E_{\ell' q'})=
\delta_{\ell\ell'}\delta_{qq'}$.
The product $E_{\ell q}^\dagger E_{\ell' q'}$ contains tensor ranks at
most $\ell+\ell'\le6$.  Anticoherence therefore implies
\[
 \langle\psi|E_{\ell q}^\dagger E_{\ell' q'}|\psi\rangle
 =\frac{\delta_{\ell\ell'}\delta_{qq'}}{17}.
\]
Thus the fifteen vectors $\sqrt{17}E_{\ell q}|\psi\rangle$ form an
orthonormal basis of $V_7$, and, with $\rho=|\psi\rangle\langle\psi|$,
\begin{equation}\label{eq:channel}
 \Phi(\rho):=\sum_{\ell=1}^{3}\sum_{q=-\ell}^{\ell}
 E_{\ell q}\rho E_{\ell q}^\dagger=\frac{I_{15}}{17}.
\end{equation}

The map $\Phi$ intertwines the spin conjugation actions.  Each tensor
rank occurs once in the input and output operator spaces, so it acts by
a scalar on each common rank.  Let
\[
 (T^{(j)}_{LM})_{m',m}=(-1)^{j-m'}\sqrt{2L+1}
 \begin{pmatrix}j&L&j\\-m'&M&m\end{pmatrix}.
\]
At rank seven the scalar is exactly
\begin{equation}\label{eq:multiplier}
 \Phi(T^{(8)}_{7M})=
 \frac{5\sqrt{69}}{8398}\,T^{(7)}_{7M}\ne0,
 \qquad M=-7,\ldots,7.
\end{equation}
This coefficient can be checked directly from the displayed $3j$
matrices.  For $M=0$, the contributions of $\ell=1,2,3$ to
$\Tr[T^{(7)\dagger}_{70}\Phi(T^{(8)}_{70})]$ are, respectively,
\[
 \frac{3\sqrt{69}}{170},\qquad
 \frac{\sqrt{69}}{306},\qquad
 -\frac{7679\sqrt{69}}{377910};
\]
their sum is $5\sqrt{69}/8398$.  The accompanying exact verifier checks
every diagonal matrix element, in addition to this scalar contraction.
Covariance gives the same multiplier for every $M$.

The right side of \eqref{eq:channel} has no rank-seven component.
Because \eqref{eq:multiplier} is nonzero, every rank-seven multipole of
$\rho$ must vanish.  Hence $|\psi\rangle$ would actually be $7$-AC.

Finally take the Hilbert--Schmidt orthonormal rectangular tensors
$F_{\ell q}:V_8\to V_{15/2}$ of ranks
$\ell=\tfrac12,\tfrac32,\tfrac52,\tfrac72$.  Their products have ranks
at most seven, so the same Gram argument would give
$2+4+6+8=20$ nonzero pairwise orthogonal images in the
sixteen-dimensional space $V_{15/2}$.  This contradiction proves the claim.
\end{proof}

\paragraph{General conditional equality test.}
For a pure $2r$-AC state with $N=(r+1)^2$, $r\ge2$, take
$E_{\ell q}:V_{N/2}\to V_{(N-2)/2}$ for $\ell=1,\ldots,r$.
There are $r(r+2)=N-1$ orthogonal images, so again the channel output
must be scalar.  With the tensor convention above its rank-$L$
multiplier is
\[
 c_{r,L}=\sum_{\ell=1}^r(-1)^{N+L+\ell+1}(2\ell+1)
 \begin{Bmatrix}
 N/2&N/2&L\\(N-2)/2&(N-2)/2&\ell
 \end{Bmatrix}.
\]
If $c_{r,2r+1}\ne0$, equality is impossible: the state becomes
$(2r+1)$-AC, whereas the one-qubit rectangular packing bound requires
$N\ge(r+1)(r+2)>(r+1)^2$.  This is a rigorous criterion for each $r$;
nonvanishing for every $r\ge2$ has not been proved here.
The verifier checks nonvanishing for $2\le r\le8$, independently
checking the physical $3j$ channel for $r=2,3$.

This argument concerns pure states at the minimal-size packing bound.
It does not assume that a Majorana constellation of an AC state is a
spherical design, nor that every optimal subspace has polyhedral symmetry.
The analogous Fisher-bound pattern for spherical designs is motivation,
not a premise of the proof.
\end{document}
